% ================= PORTADA =================
\thispagestyle{empty}
\begin{tikzpicture}[remember picture,overlay]
  \foreach \x in {0,5,...,150} \draw[borde!60,line width=0.2pt] ([xshift=\x mm]current page.south west) -- ([xshift=\x mm]current page.north west);
  \foreach \y in {0,5,...,215} \draw[borde!60,line width=0.2pt] ([yshift=\y mm]current page.south west) -- ([yshift=\y mm]current page.south east);
\end{tikzpicture}
\vspace*{18mm}
{\ttfamily\small\color{suave} CDIV · IMERL · FING · 2S 2026}\par\vspace{3mm}
{\plexbold\fontsize{26}{30}\selectfont Cálculo 1}\par\vspace{1mm}
{\plexbold\fontsize{17}{21}\selectfont\color{rojo} Prácticos resueltos}\par\vspace{4mm}
{\large Todos los ejercicios que entran al primer parcial,\\ uno por uno, paso a paso.}\par\vspace{10mm}
\begin{sello}{CÓMO ESTÁ ARMADO}
\textbf{Parte 0:} todas las propiedades que usás en el curso, en orden, de la más simple (la conmutativa) a la más compleja (Weierstrass). Cada una tiene un número \textbf{P1, P2, \dots}\\[3pt]
\textbf{Capítulos 1 a 5:} los prácticos oficiales resueltos. En cada paso, una cajita gris \textbf{TEORÍA QUE USO} te dice qué propiedad se aplica, con su número. Si no te acordás, vas a la Parte 0.
\end{sello}
\vfill
{\ttfamily\scriptsize\color{suave} Hecho sobre los prácticos oficiales 1 a 5 y tus apuntes de clase. Todas las cuentas verificadas por computadora (sympy).}
\newpage

{\hypersetup{linkcolor=tinta}\renewcommand{\contentsname}{Índice}\tableofcontents}
\newpage

\chapter*{Cómo usar este cuaderno}
\markboth{CÓMO USAR ESTE CUADERNO}{}

\begin{memo}{LAS CAJAS}
\textbf{\color{rojo}Roja doble:} definición o teorema que hay que saber literal.\\
\textbf{\color{acero}Gris fina, «TEORÍA QUE USO»:} la propiedad que se aplica en ese paso, con su número de la Parte 0.\\
\textbf{\color{acero}Gris con borde:} método, el procedimiento paso a paso.\\
\textbf{\color{verde}Verde:} ayuda-memoria. \textbf{\color{teal}Teal:} intuición o analogía.\\
\textbf{\color{rojo}Barra roja:} trampa típica. \textbf{\color{naranja}Barra naranja:} el ejercicio depende de una figura del EVA que no tengo; te dejo el método para que lo termines mirando la figura.
\end{memo}

\section*{Qué entra y qué no}
Este cuaderno cubre las secciones de los prácticos que el cronograma asigna antes del primer parcial:

\begin{center}\small
\begin{tabularx}{\linewidth}{@{}l X@{}}
\toprule
\textbf{Capítulo} & \textbf{Secciones resueltas} \\ \midrule
1 · Conjuntos y funciones & 1.1 Conjuntos · 1.2 Funciones · 1.3 Modelado \\
2 · Número real & 2.1 Axiomas · 2.2 Funciones reales · 2.3 Completitud · 2.4 Funciones y completitud \\
3 · Integración & 3.1 a 3.7 (introducción, acotaciones, Darboux, integrabilidad, cambio de variable, cálculo, aplicaciones) \\
4 · Límites & 4.1 a 4.5 (geométrica, definición, cálculo, infinito, continuidad) \\
5 · Funciones continuas & 5.1 Continuidad · 5.2 Bolzano · 5.3 Weierstrass \\ \bottomrule
\end{tabularx}
\end{center}

Complementarios, Talleres y Laboratorio \textbf{no} están: el cronograma no los asigna para el parcial. Si te sobra tiempo, el libro del parcial tiene los ejercicios con más probabilidad.

\begin{tip}{CÓMO LEER CADA EJERCICIO}
Primera pasada: leelo entero. Segunda: tapá la resolución, leé solo el enunciado y escribí los pasos vos. Cuando te trabes, mirá solo la cajita \textbf{TEORÍA QUE USO} del paso siguiente: te dice qué herramienta usar sin darte la cuenta.
\end{tip}
